\clearpage\pagestyle{matSimple}\phantomsection\label{sources}\pdfbookmark[0]{Sources}{bm-sources}
\matPageTitre{Sources et édition}{Références officielles · France}
{\fontsize{9.5}{13.5}\selectfont \textbf{1. Programme d’enseignement de l’école maternelle}\par Arrêté publié au BO n° 41 du 31 octobre 2024, applicable à la rentrée 2025 ; version consolidée intégrant le BO n° 19 du 7 mai 2026. Domaine « L’acquisition des premiers outils mathématiques » et partie « Se repérer dans l’espace ».\par\href{https://eduscol.education.gouv.fr/sites/default/files/document/programme-cycle-1-consolide-127565.pdf}{\textcolor{otGreenDark}{Consulter la source officielle}}\par\medskip
\textbf{2. Ressources d’accompagnement Éduscol}\par Ressources pour le cycle 1 : construction du nombre, résolution de problèmes, formes et grandeurs, motifs organisés.\par\href{https://eduscol.education.gouv.fr/}{\textcolor{otGreenDark}{Consulter la source officielle}}\par\medskip
\textbf{3. Choix faits dans ce cahier}\par Les quantités à dénombrer ne dépassent pas six ; les comparaisons globales portent sur des collections plus grandes, dont l’une a au moins deux fois plus d’objets que l’autre. La comptine est récitée jusqu’à douze, mais les chiffres écrits ne dépassent pas six. Les solides ne sont pas dessinés en perspective : on travaille leurs empreintes, avec de vrais solides. Les assemblages comportent au plus cinq pièces. Les doigts sont montrés en commençant par le pouce, mais toute autre façon de montrer un nombre est juste. Chaque fiche commence par une activité avec de vrais objets, décrite dans le guide de l’adulte.\par}
